\chapter{Surveillance et réentraînement}\label{ch:ml:monitoring-and-retraining}

Un mardi, le modèle a cessé de trader~: chaque prédiction se trouvait à l'intérieur du seuil. Rien n'avait échoué. Un fournisseur en amont avait
commencé à envoyer des volumes en lots plutôt qu'en actions, et le modèle les redimensionnait docilement. Aucun test n'avait cassé, car chaque
nombre restait un nombre~; le seul signe était que le desk n'avait rien à faire. Ce chapitre construit les moniteurs qui auraient
averti quelqu'un avant le déjeuner, et interrogent chacun sur la question qui décide de son utilité~: combien de temps faut-il pour
détecter une vraie défaillance lorsqu'il n'est autorisé à crier au loup qu'une fois par mois~? Trois défaillances sont introduites dans une simulation de production~: un
changement d'unité en amont, un renversement de l'effet de la variable la plus forte, et une variable dont l'effet s'estompe pendant que sa dispersion augmente.
Les moniteurs des entrées voient la première le jour même où elle survient et sont aveugles au deuxième~; le moniteur des résultats voit le deuxième en un
jour et ne voit pas la première avant des mois~; rien ne repère la perte de précision de la troisième, et les moniteurs des entrées ne la détectent que par son
effet secondaire, au bout de six semaines.

\section{Que surveiller~: entrées, sorties, résultats}

\begin{definition}[Surveillance des modèles, dérive des prédictions]\label{def:ml:monitoring-and-retraining:monitoring}
La \emph{surveillance des modèles}\index{surveillance des modèles} est la comparaison continue, pour un modèle en production, des entrées, des sorties et
des résultats réalisés avec ce qui était attendu lors de sa validation, et des alarmes qui sollicitent une personne lorsque les écarts dépassent
ce que le hasard permet. La \emph{dérive des prédictions}\index{dérive des prédictions} est un changement de la distribution des sorties du modèle par rapport
à la période de validation, quelle qu'en soit la cause.
\end{definition}

Le chapitre~12 distingue le \omterm{def:ml:online-learning-and-drift:drift}{décalage des covariables} (covariate shift), qui touche les entrées, de la \omterm{def:ml:online-learning-and-drift:drift}{dérive conceptuelle} (concept drift), qui touche la manière dont les entrées se rapportent à la
cible, et construit des \omterm{def:ml:online-learning-and-drift:detector}{détecteurs de dérive} fondés sur les erreurs d'un modèle. Un modèle en production est surveillé à trois niveaux, car chacun repère des défaillances
que les autres manquent. Les entrées sont disponibles immédiatement et en masse~: chaque variable de chaque valeur, chaque jour, comparée à la
référence d'entraînement, après les contrôles de schéma et de livraison du chapitre~15 et les contrôles de fraîcheur du chapitre~24. Les sorties sont accessibles
tout aussi rapidement~: la distribution des prédictions, la part des valeurs dont la prévision franchit le seuil de trading, les positions et
le taux de rotation qu'elles impliquent. Les résultats (le coefficient d'information réalisé, le P\&L) n'arrivent qu'après l'horizon de prévision, et ils
sont bruités~: un modèle dont le coefficient d'information quotidien vrai vaut 0.10 présente un écart-type d'environ 0.07 d'un jour à
l'autre. Ce sont les résultats qui comptent, et ce sont eux qui parlent le plus tard.

Le modèle de production de ce chapitre est volontairement simple, afin que chaque défaillance soit connue exactement. Chaque jour, 200 valeurs portent cinq
variables standardisées~; les rendements du lendemain dépendent des trois premières avec des coefficients 0.08, $-0.06$ et 0.05, plus un bruit de
variance unitaire. Un modèle linéaire est ajusté sur 60 jours, puis tourne 400 jours de plus~; il trade les valeurs dont la prévision est supérieure en
valeur absolue à la médiane d'entraînement, soit la moitié d'entre elles. (Le modèle de carnet d'ordres du chapitre~29 est surveillé avec le même code en fin de
livre.) Trois défaillances commencent au jour 200, chacune dans 20 simulations indépendantes (\cref{lst:ml:mon:day})~:
\begin{itemize}
\item \emph{changement d'unité}~: la deuxième variable arrive cent fois plus petite, comme le feraient des volumes en lots~;
\item \emph{renversement}~: l'effet de la première variable change de signe, un changement de régime que les entrées ne montrent pas~;
\item \emph{érosion lente}~: sur 150 jours, l'effet de la troisième variable tombe à zéro pendant que sa dispersion augmente de moitié.
\end{itemize}

\section{Statistiques de dérive et conception des alarmes}

\begin{definition}[Indice de stabilité de la population]\label{def:ml:monitoring-and-retraining:psi}
On découpe un échantillon de référence d'une variable en $B$ classes (ici ses déciles) de parts $r_1,\dots,r_B$, et notons $q_1,\dots,q_B$ les
parts d'un nouvel échantillon dans les mêmes classes. L'\emph{indice de stabilité de la population}\index{indice de stabilité de la population} (population stability index, PSI) est
\[
\mathrm{PSI}=\sum_{b=1}^{B}(q_b-r_b)\ln\frac{q_b}{r_b},
\]
la divergence de Kullback--Leibler symétrisée entre les deux distributions regroupées en classes (les classes vides sont ramenées à une petite part).
\end{definition}

Cet indice vient du scoring de crédit, où les règles empiriques veulent qu'en dessous de 0.10 une population ait peu changé et qu'au-dessus de 0.25
elle a nettement changé~; Yurdakul et Naranjo soulignent que ces repères sont utilisés sans tenir compte de leurs
taux d'erreur. La statistique de Kolmogorov--Smirnov du livre~4 (chapitre~12), le plus grand écart entre les deux fonctions de répartition
empiriques, est l'autre outil de référence (\cref{lst:ml:mon:psi}). Les deux sont calculées chaque jour sur chaque variable par rapport à la
référence d'entraînement, et le moniteur retient la plus grande sur les cinq variables. Avec beaucoup de variables, un test par variable multiplie les
fausses alarmes~; Rabanser et ses coauteurs, en comparant les façons de détecter un décalage de jeu de données, ont trouvé que les tests à deux échantillons portant sur une version réduite
des entrées étaient les plus efficaces~: un test multivarié unique, comme le \omterm{def:ml:generative-models-and-synthetic-data:tests}{test à deux échantillons par classifieur} du chapitre~16. Les moniteurs de sorties calculent l'indice sur les prédictions
et suivent la part des valeurs à l'intérieur du seuil de trading, c'est-à-dire là où la défaillance du mardi se manifesterait en premier.

Une alarme est un seuil, et un seuil est un taux de fausses alarmes. Le chapitre fixe une fausse alerte par mois, soit une journée de séance sur 21,
et calibre chaque seuil sur 20 simulations de 400 jours sans défaillance (\cref{lst:ml:mon:calib}). Une alerte est comptée lorsqu'une alarme
démarre, et non pour chaque jour où elle dure, car une statistique qui reste au-dessus de son seuil pendant une semaine ne sollicite la personne d'astreinte qu'une seule fois.
Les seuils calibrés valent 0.115 pour l'indice maximal sur les variables, 0.115 pour la plus grande statistique de Kolmogorov--Smirnov, 0.089
pour l'indice des prédictions et 0.07 pour la variation de la part située à l'intérieur du seuil. Le premier dépasse la règle empirique du
scoring de crédit~: avec 200 valeurs par jour, chaque décile en contient 20, et le seul bruit d'échantillonnage porte le plus grand des cinq indices au-delà de
0.10 plus d'un jour sur vingt. Un seuil recopié d'une règle empirique alerterait plus souvent que ce dont l'équipe était convenue, et
une équipe trop souvent alertée cesse de lire les alertes.

Deux nombres décrivent la réaction de chaque moniteur à une défaillance. La \emph{première alarme} est le nombre de jours écoulés entre la défaillance et la
première alerte du moniteur~; comme chaque moniteur émet une fausse alerte par mois, un moniteur qui ne voit rien alerte encore après environ deux
semaines, et une première alarme n'est une détection que si elle survient bien avant. La \emph{alarme persistante} est le nombre de jours
jusqu'à la première fenêtre glissante de 21 jours comptant au moins cinq jours d'alarme, une alerte qui revient sans cesse~; sans défaillance, la médiane
d'aucun moniteur n'atteint une telle fenêtre dans les 200 jours.

\section{Alarmes de performance sur des résultats bruités}

\begin{definition}[Alarme de performance]\label{def:ml:monitoring-and-retraining:performance-alarm}
Une \emph{alarme de performance}\index{alarme de performance} est un moniteur portant sur les résultats réalisés (le coefficient d'information, le taux de réussite (hit rate),
le P\&L ou les mark-outs des transactions du modèle) qui se déclenche lorsque leur niveau tombe sous ce que la validation avait conduit l'établissement à attendre, d'un écart plus grand
que ne le permet le bruit propre aux résultats au taux de fausses alarmes choisi.
\end{definition}

L'\omterm{def:ml:monitoring-and-retraining:performance-alarm}{alarme de performance} du chapitre est le test CUSUM unilatéral du livre~7 (chapitre~13) appliqué au coefficient d'information quotidien
(\cref{lst:ml:mon:cusum})~: il cumule le déficit du coefficient de chaque jour sous son niveau attendu de 0.103, moins une
marge $k$ égale à la moitié de son écart-type quotidien (0.035), et déclenche lorsque la somme dépasse $h$, puis repart de zéro. Pour une fausse
alerte par mois, $h$ vaut 0.11, soit environ un jour et demi d'écart-type~: le moniteur peut se permettre d'attendre très peu.

Cela indique ce qu'il peut voir. Une baisse du coefficient moyen de $\Delta$, mesurée avec un bruit quotidien $\sigma$, il faut environ
\[
n\approx\Bigl(\frac{(z_{0.95}+z_{0.90})\,\sigma}{\Delta}\Bigr)^2
\]
jours pour être distingué du bruit avec un test unilatéral au niveau de 5~\% et une puissance de 90~\%. Après le renversement, le coefficient passe de 0.103 à
0.003, et $n$ vaut environ 4 jours~; après le changement d'unité, il tombe à 0.087 et $n$ vaut environ 150 jours~; à la fin de l'érosion lente, il
est tombé à 0.081, et $n$ vaut environ 80 jours, mais l'érosion met 150 jours à y parvenir. Les résultats ne voient pas de petites pertes de
précision sur aucun horizon qu'un desk accepterait~; ils repèrent les grandes rapidement.

\begin{table}[htbp]
\centering\small
\small\setlength{\tabcolsep}{4pt}%
\begin{tabular}{@{}>{\raggedright\arraybackslash}p{3.0cm}ccccccc@{}}
\toprule
 & aucune & \multicolumn{2}{c}{changement d'unité} & \multicolumn{2}{c}{renversement} & \multicolumn{2}{c}{érosion lente}\\
\cmidrule(lr){2-2}\cmidrule(lr){3-4}\cmidrule(lr){5-6}\cmidrule(lr){7-8}
moniteur & première & première & persistante & première & persistante & première & persistante\\
\midrule
PSI maximal sur les variables & 12 & 0 & 0 & 12 & jamais & 12 & 42\\
KS maximal sur les variables & 10 & 0 & 0 & 10 & jamais & 8 & 60\\
PSI des prédictions & 17 & 0 & 0 & 17 & jamais & 19.5 & jamais\\
part à l'intérieur du seuil & 11 & 0 & 0 & 11 & jamais & 9.5 & jamais\\
CUSUM sur l'IC & 13 & 9.5 & jamais & 1 & 0 & 13 & jamais\\
\bottomrule
\end{tabular}
\caption{Jours écoulés depuis la défaillance (jour 200) jusqu'à la première alarme de chaque moniteur et jusqu'à une alarme persistante (cinq jours d'alarme sur 21), médianes
sur 20 simulations, chaque moniteur étant calibré pour une fausse alerte par mois~; «~aucune~» désigne le même décompte sur les simulations sans défaillance, soit
le score d'un moniteur aveugle~; «~jamais~» signifie l'absence d'alarme persistante dans les 200 jours qui suivent. Données~: \texttt{ml\_monitor.delays}.}
\label{tab:ml:mon:delays}
\end{table}

Le \Cref{tab:ml:mon:delays} présente le résultat du chapitre. Le changement d'unité est repéré par chaque moniteur d'entrées et de sorties dès son premier jour~:
le PSI de la deuxième variable vaut environ 7, soit soixante fois son seuil, et la part des valeurs situées à l'intérieur du seuil de négociation passe de
50~\% à 63~\%, soit un quart de valeurs négociées en moins. Le moniteur de résultats ne déclenche qu'une alerte, après une médiane de 9.5 jours, et n'insiste jamais. Quant à
l'inversion est invisible à chaque moniteur d'entrées et de sorties (leurs colonnes sont égales à la colonne «~aucune~»~: les entrées et la
distribution des prédictions n'ont pas changé du tout), et le CUSUM déclenche une alerte en deux jours dans 13 simulations sur 20 et continue d'en déclencher. Enfin,
la perte de précision de la dégradation lente n'est vue par rien~; sa dispersion croissante est repérée par les moniteurs d'entrées, en alarme persistante après une médiane de 42
jours pour le PSI et 60 pour la statistique de Kolmogorov--Smirnov (la \cref{fig:ml:mon:fade}). Si la dégradation était survenue sans le changement
de dispersion, rien dans le tableau ne l'aurait détectée.

\begin{omfigure}
\begin{tikzpicture}
\begin{axis}[width=11.5cm,height=6cm,xlabel={jour de production},ylabel={jours d'alarme par mois},xmin=0,xmax=400,ymin=0,ymax=22,
  tick label style={font=\small},label style={font=\small},grid=major,grid style={gray!25},legend pos=north west,
  legend cell align=left,legend style={font=\small},table/col sep=comma]
\addplot[omProp,very thick,mark=*,mark size=1.3pt] table[x=day,y=psi]{figdata/ml/27-monitoring-and-retraining/fade.csv};
\addplot[omDef,very thick,dashed,mark=square*,mark size=1.3pt,mark options={solid}] table[x=day,y=ks]{figdata/ml/27-monitoring-and-retraining/fade.csv};
\addplot[gray,thick,mark=triangle*,mark size=1.5pt] table[x=day,y=predpsi]{figdata/ml/27-monitoring-and-retraining/fade.csv};
\addplot[omThm,very thick,dotted,mark=diamond*,mark size=1.6pt,mark options={solid}] table[x=day,y=cusum]{figdata/ml/27-monitoring-and-retraining/fade.csv};
\addplot[black,thin,dashed,forget plot] coordinates {(200,0) (200,22)};
\legend{PSI de la variable la plus forte,KS de la variable la plus forte,PSI des prédictions,CUSUM sur l'IC}
\end{axis}
\end{tikzpicture}

\omcaption{Jours d'alarme par mois (21 jours) sous la dégradation lente, sur des fenêtres de 25 jours, moyennés sur 20 simulations~; la dégradation commence au jour
200 (en tirets) et se termine au jour 350. Un moniteur aveugle déclenche une alarme environ un jour par mois. Données~:
\texttt{ml\_monitor.alarm\_rates}.}\label{fig:ml:mon:fade}
\end{omfigure}

\begin{method}[Conception des moniteurs d'un modèle]\label{met:ml:mon:design}
\begin{enumerate}
\item Surveiller les entrées, les sorties et les résultats~: chaque couche détecte des défaillances que les autres ne peuvent pas voir.
\item Choisir le taux de fausses alertes auquel l'équipe répondra effectivement, et calibrer chaque seuil en fonction de ce taux sur l'historique ou par simulation
sans défaillance~; ne pas reprendre les seuils des règles empiriques.
\item Compter les alertes au début d'une alarme, et distinguer une première alerte d'une alarme persistante.
\item Mesurer le délai de chaque moniteur sur des défaillances introduites, et le comparer à celui d'un moniteur aveugle~; un moniteur qui ne bat jamais
cette colonne n'est que du bruit.
\item Admettre que les moniteurs de résultats ne voient rapidement que les grosses pertes~; couvrir les petites par une revalidation programmée.
\end{enumerate}
\end{method}

\section{Déploiements fantôme, canari et champion--challenger}

\begin{definition}[Déploiement fantôme, champion--challenger]\label{def:ml:monitoring-and-retraining:shadow}
Dans un \emph{déploiement fantôme}\index{déploiement fantôme} (shadow deployment), un nouveau modèle reçoit les entrées de production et enregistre ses prédictions,
mais ses décisions ne sont pas exécutées~; ses résultats sont calculés comme s'ils l'avaient été. Le \emph{champion--challenger}\index{champion--challenger}
est la discipline qui consiste à garder en place le modèle de production (le champion) jusqu'à ce qu'un challenger, exécuté à ses côtés sur les mêmes données,
ait montré, à l'aide d'un test choisi à l'avance, qu'il est meilleur.
\end{definition}

Une exécution en \omterm{def:ml:monitoring-and-retraining:shadow}{déploiement fantôme} ne coûte que du calcul et fournit une comparaison nette, car le champion et le challenger voient les mêmes valeurs
aux mêmes jours. Après l'inversion, le CUSUM déclenche une alerte en quelques jours. Un challenger est réajusté sur les 40 jours
qui suivent (son coefficient sur la première variable vaut $-0.081$ en moyenne sur les simulations, contre la vraie valeur $-0.08$) et fonctionne en \omterm{def:ml:monitoring-and-retraining:shadow}{déploiement fantôme}
à partir du jour 240. Son IC moyen vaut 0.105, celui du champion 0.003. La comparaison repose sur le test séquentiel toujours valide du livre~7
(chapitre~21) sur les différences quotidiennes (le \cref{lst:ml:mon:shadow}), qu'on peut examiner chaque jour sans
gonfler son taux d'erreur~; l'erreur type est estimée à partir des jours observés, si bien que le test ne donne aucun verdict pendant dix jours, car sur
trois jours, un petit écart type d'échantillon rend toute différence certaine. Le challenger est promu dans les 20 simulations, après
une médiane de 10 jours en \omterm{def:ml:monitoring-and-retraining:shadow}{déploiement fantôme} et au plus 19 jours (la \cref{fig:ml:mon:shadow}). Sans défaillance, un challenger réajusté sur 40 jours
ordinaires est promu dans 2 simulations sur 100, dans la limite des 5~\% du test.

\begin{omfigure}
\begin{tikzpicture}
\begin{axis}[width=11cm,height=5.8cm,xlabel={jour de déploiement fantôme},ylabel={valeur p toujours valide},ymode=log,xmin=1,xmax=20,
  ymin=1e-7,ymax=2,tick label style={font=\small},label style={font=\small},grid=major,grid style={gray!25},
  table/col sep=comma]
\addplot[omProp,thick,mark=*,mark size=1.2pt] table[x=day,y=run0]{figdata/ml/27-monitoring-and-retraining/shadow.csv};
\addplot[omDef,thick,mark=square*,mark size=1.2pt] table[x=day,y=run1]{figdata/ml/27-monitoring-and-retraining/shadow.csv};
\addplot[omThm,thick,mark=triangle*,mark size=1.4pt] table[x=day,y=run2]{figdata/ml/27-monitoring-and-retraining/shadow.csv};
\addplot[gray,thick,mark=diamond*,mark size=1.4pt] table[x=day,y=run3]{figdata/ml/27-monitoring-and-retraining/shadow.csv};
\addplot[black,thick,mark=o,mark size=1.2pt] table[x=day,y=run4]{figdata/ml/27-monitoring-and-retraining/shadow.csv};
\addplot[black,thin,dashed] coordinates {(1,0.05) (20,0.05)};
\end{axis}
\end{tikzpicture}

\omcaption{Valeurs p toujours valides de l'avantage du challenger sur le champion, jour par jour en \omterm{def:ml:monitoring-and-retraining:shadow}{déploiement fantôme} après l'inversion, pour cinq
des 20 simulations~; le test ne donne aucun verdict pendant dix jours~; en tirets~: 0.05, seuil auquel le challenger est promu. Données~:
\texttt{ml\_monitor.shadow\_run}.}\label{fig:ml:mon:shadow}
\end{omfigure}

Les comparaisons en \omterm{def:ml:monitoring-and-retraining:shadow}{déploiement fantôme} mesurent des prévisions, non la réaction du marché à leur exploitation~; un modèle dont les ordres font bouger les cours est comparé
dans un déploiement canari (livre~7, chapitre~21), en négociant une petite part du capital. L'intérêt du déploiement canari tient à ce qu'il limite. Le
chapitre promeut un challenger affecté d'une erreur de signe sur sa deuxième variable, un bogue qu'un \omterm{def:ml:monitoring-and-retraining:shadow}{déploiement fantôme} aurait détecté, mais supposons qu'il ne l'ait pas repéré~:
le CUSUM le fait revenir en arrière après une médiane de 3 jours (au plus 18). L'IC sacrifié, cumulé sur ces jours
et pondéré par le capital, vaut 0.26 IC-jours quand le challenger négocie tout le portefeuille et 0.026 quand il n'en négocie qu'un dixième. Le
retour en arrière est lui-même une action de registre (le \cref{lst:ml:mon:rollback})~: la version qui était en production auparavant est de nouveau promue,
avec sa raison, dans la \omterm{def:ml:experiment-tracking-and-reproducibility:registry}{piste d'audit} du chapitre~25.

\section{Réentraîner, revenir en arrière ou retirer}

Chaque alarme appelle une action différente, et la mauvaise cache la défaillance. Une alarme sur les entrées sans baisse des résultats (le changement
d'unité) appelle la correction des données~: réentraîner sur des volumes par lots apprendrait au modèle une nouvelle échelle pour une variable et enfouirait
l'erreur du fournisseur dans ses coefficients, jusqu'à ce que le fournisseur la corrige. Une alarme sur les résultats, avec des entrées calmes (l'inversion), signifie que la
relation a changé~: réduire la taille du modèle ou l'arrêter (le kill switch du livre~11, chapitre~27, pour un modèle qui négocie à vitesse élevée),
réajuster sur les données postérieures au changement, dès qu'il y en a assez, et le promouvoir par \omterm{def:ml:monitoring-and-retraining:shadow}{déploiement fantôme} puis par déploiement canari. Une perte lente qu'aucun moniteur
ne voit appelle un réentraînement programmé (le \omterm{def:ml:online-learning-and-drift:schedule}{calendrier de réentraînement} du chapitre~12) et une revalidation périodique du modèle par rapport à un
indice de référence simple. Un modèle est retiré lorsque ses challengers réajustés ne battent plus un modèle nul~: quand l'effet a disparu, réentraîner
ne fait qu'ajuster du bruit plus vite. Sculley et ses coauteurs rangent les «~changements du monde extérieur~» parmi les dettes des
systèmes d'apprentissage automatique, et Breck et ses coauteurs font de la surveillance la moitié de leur grille de préparation d'un modèle à la production. Dans les entretiens
d'ingénieurs en apprentissage automatique menés par Shankar et ses coauteurs, les fausses alertes constituaient le point de friction le plus souvent évoqué, et
la fatigue d'alerte qui s'ensuit conduit les équipes à faire taire les alertes et à manquer les vraies baisses~: c'est la raison de compter les alertes.

\begin{method}[Après une alarme]\label{met:ml:mon:after}
\begin{enumerate}
\item Alarme sur les entrées ou les sorties~: vérifier d'abord le chemin des données (unités, schéma, changements de fournisseur)~; corriger les données, et non le modèle.
\item Alarme sur les résultats avec des entrées calmes~: réduire ou arrêter la négociation, réajuster sur les données postérieures au changement, et ne promouvoir que par \omterm{def:ml:monitoring-and-retraining:shadow}{déploiement fantôme} et
déploiement canari.
\item Un challenger est promu grâce à un test séquentiel choisi à l'avance, et ramené en arrière par l'\omterm{def:ml:monitoring-and-retraining:performance-alarm}{alarme de performance}, deux opérations consignées dans
le registre.
\item Réentraîner selon un calendrier pour les pertes qu'aucun moniteur ne voit~; retirer le modèle lorsque le réentraînement ne bat plus un modèle nul.
\end{enumerate}
\end{method}

\section{Tutoriel~: le mardi où le modèle s'est tu}\label{tut:ml:monitoring-and-retraining:tut}

\begin{tutorial}
\textbf{Objectif.} Injecter trois défaillances dans une simulation de production, calibrer cinq moniteurs à une fausse alerte par mois, mesurer
leurs délais, promouvoir un challenger issu du \omterm{def:ml:monitoring-and-retraining:shadow}{déploiement fantôme} et ramener en arrière un challenger défaillant. \textbf{État final~:} le \cref{tab:ml:mon:delays},
la \cref{fig:ml:mon:fade}, la \cref{fig:ml:mon:shadow}.
\begin{tutsteps}
\item \textbf{Les trois défaillances.}
\omcode{code/ml/27-monitoring-and-retraining/python/ml_monitor.py}{44}{58}{Une journée de données de production, avec les défaillances injectées.}\label{lst:ml:mon:day}
\item \textbf{Statistiques de dérive.}
\omcode{code/firm/mlmonitor/firm_mlmonitor.py}{49}{65}{L'indice de stabilité de la population et la statistique de Kolmogorov--Smirnov par rapport à une référence préparée.}\label{lst:ml:mon:psi}
\item \textbf{Seuils pour un taux d'alertes donné.}
\omcode{code/firm/mlmonitor/firm_mlmonitor.py}{94}{110}{Débuts d'alarme, et seuil qui déclenche une alerte à un taux choisi.}\label{lst:ml:mon:calib}
\item \textbf{L'\omterm{def:ml:monitoring-and-retraining:performance-alarm}{alarme de performance}.}
\omcode{code/firm/mlmonitor/firm_mlmonitor.py}{75}{86}{Un CUSUM unilatéral sur le coefficient d'information quotidien.}\label{lst:ml:mon:cusum}
\item \textbf{Le test sur \omterm{def:ml:monitoring-and-retraining:shadow}{déploiement fantôme}.}
\omcode{code/firm/mlmonitor/firm_mlmonitor.py}{118}{127}{Valeurs p toujours valides de l'avantage d'un challenger, sans verdict avant dix jours.}\label{lst:ml:mon:shadow}
\item \textbf{Retour en arrière.}
\omcode{code/firm/mlmonitor/firm_mlmonitor.py}{130}{143}{Remettre la version de production précédente en place via le registre.}\label{lst:ml:mon:rollback}
\item \textbf{Exécuter} \texttt{ml\_monitor.thresholds()}, \texttt{delays()}, \texttt{shadow()}, \texttt{rollbacks()} et
\texttt{fig\_monitor.py} (environ 15 secondes sur un cœur).
\end{tutsteps}
\textbf{Pour aller plus loin.} Regrouper les 20 derniers jours de chaque variable en un seul indice et le calibrer de la même façon~: il repère plus tôt
la dispersion de la lente érosion, mais ses fausses alarmes arrivent par salves, et une alerte par mois n'a pas le même sens pour lui
(l'\cref{exo:ml:monitoring-and-retraining:7}).
\end{tutorial}

\section{Construction~: surveillance des modèles}\label{bld:ml:monitoring-and-retraining:mlmonitor}

\begin{build}
\bfield{Rôle} Moniteurs sur les entrées, les sorties et les résultats d'un modèle de production, calibrés à un taux de fausses alertes, et les
décisions de promotion et de retour en arrière qu'ils alimentent.
\bfield{Interface} \texttt{Reference(ref, bins)}~; \texttt{psi(ref, x)}~; \texttt{ks(ref, x)}~; \texttt{daily\_ic(pred, y)}~;
\texttt{Cusum(target, k, h).update(v)}~; \texttt{calibrate\_daily(null, rate)}~; \texttt{onsets(alarms)}~;
\texttt{calibrate\_pages(null, rate)}~; \texttt{first\_alarm(stats, threshold, start)}~;
\texttt{shadow\_verdict(ic\_challenger, ic\_champion, tau, min\_days)} sur \texttt{firm.abtest}~; \texttt{rollback(registry, name, ts, reason)} sur \texttt{firm.exptrack}.
\bfield{Règles} Chaque seuil est calibré sur des données sans défaillance à un taux énoncé~; les alertes sont comptées aux débuts d'alarme~; les promotions
utilisent un test fixé à l'avance~; les retours en arrière passent par le registre, avec un motif.
\bfield{Tests d'acceptation} \texttt{code/firm/mlmonitor/tests/}~: l'indice est nul sur sa propre référence et croît avec un décalage~; la
statistique de Kolmogorov--Smirnov et le coefficient d'information concordent avec SciPy, ex aequo compris~; le CUSUM déclenche une alerte sur une baisse puis se réinitialise~;
les seuils calibrés déclenchent des alertes à leur taux~; le test sur \omterm{def:ml:monitoring-and-retraining:shadow}{déploiement fantôme} ne rend aucun verdict avant dix jours, et ne promeut un challenger sans avantage que dans
au plus 5~\% des simulations~; un retour en arrière rétablit la version précédente et préserve intacte la \omterm{def:ml:experiment-tracking-and-reproducibility:registry}{piste d'audit}.
\bfield{Pour aller plus loin} Dérive multivariée (un \omterm{def:ml:generative-models-and-synthetic-data:tests}{test à deux échantillons par classifieur}, chapitre~16) et moniteurs sur le P\&L et les mark-out, calibrés de la
même façon.
\end{build}

\begin{omsources}
\item S.~Rabanser, S.~Günnemann et Z.~C.~Lipton, ``Failing loudly: an empirical study of methods for detecting dataset
shift'', \emph{NeurIPS}, 2019.
\item E.~Breck, S.~Cai, E.~Nielsen, M.~Salib et D.~Sculley, ``The ML test score: a rubric for ML production readiness and
technical debt reduction'', \emph{IEEE International Conference on Big Data}, 2017.
\item B.~Yurdakul et J.~D.~Naranjo, ``Statistical properties of the population stability index'', \emph{Journal of Risk Model
Validation}, 2020.
\item D.~Sculley et coauteurs, ``Hidden technical debt in machine learning systems'', \emph{NeurIPS}, 2015.
\item S.~Shankar, R.~Garcia, J.~M.~Hellerstein et A.~G.~Parameswaran, ``Operationalizing machine learning: an interview
study'', arXiv:2209.09125, 2022~; publié dans \emph{Proceedings of the ACM on Human-Computer Interaction}, 2024.
\end{omsources}

\section{Exercices}

\begin{exercise}[$\star$]\label{exo:ml:monitoring-and-retraining:1}
Pour chacune des trois couches (entrées, sorties, résultats), donner une défaillance qu'elle voit et que les deux autres ne voient pas, parmi
les trois du chapitre.
\end{exercise}

\begin{exercise}[$\star$]\label{exo:ml:monitoring-and-retraining:2}
Expliquer pourquoi le seuil calibré sur le plus grand indice de variable (0.115) dépasse la règle empirique de 0.10 du scoring
de crédit, et ce que donnerait un seuil de 0.10.
\end{exercise}

\begin{exercise}[$\star$]\label{exo:ml:monitoring-and-retraining:3}
Expliquer pourquoi un moniteur aveugle, calibré à une fausse alerte par mois, signale tout de même une première alarme environ deux semaines après une défaillance.
\end{exercise}

\begin{exercise}[$\star\star$]\label{exo:ml:monitoring-and-retraining:4}
Vérifier le délai d'environ 150 jours que la figure du chapitre donne au moniteur de résultats pour voir le changement d'unité, à partir des coefficients d'information
0.103 et 0.087 et de l'écart-type quotidien 0.071. Calculer le nombre de jours qu'il faudrait à une firme négociant 2\,000 valeurs.
\end{exercise}

\begin{exercise}[$\star\star$]\label{exo:ml:monitoring-and-retraining:5}
Expliquer pourquoi le test sur \omterm{def:ml:monitoring-and-retraining:shadow}{déploiement fantôme} ne rend aucun verdict avant dix jours, et ce qui se passerait sans cette règle.
\end{exercise}

\begin{exercise}[$\star\star$]\label{exo:ml:monitoring-and-retraining:6}
\emph{Trouver la faille.} «~L'alarme PSI s'est déclenchée après le changement de fournisseur, donc on a réentraîné le modèle sur les nouvelles données, et l'alarme
s'est arrêtée.~»
\end{exercise}

\begin{exercise}[$\star\star\star$]\label{exo:ml:monitoring-and-retraining:7}
\emph{Programmation.} Ajouter le moniteur regroupé de la dernière étape du tutoriel (le plus grand indice de variable sur les 20 derniers jours d'échantillons),
calibrer à une alerte par mois, et donner ses délais. Expliquer pourquoi ses fausses alarmes arrivent par salves.
\end{exercise}

\begin{exercise}[$\star\star\star$]\label{exo:ml:monitoring-and-retraining:8}
Concevoir les moniteurs pour une érosion sans changement de dispersion~: l'effet de la troisième variable tombe à zéro en 150 jours et sa
distribution ne change pas. Indiquer quelle couche peut le voir, à quelle vitesse, et ce qu'on ferait à la place.
\end{exercise}

\section{Problème~: le mardi où le modèle s'est tu}

\begin{problem}[{Problème du week-end --- le mardi où le modèle s'est tu}]\label{pb:ml:monitoring-and-retraining:1}
La simulation de production du chapitre, trois défaillances injectées, cinq moniteurs.

\medskip\noindent\textbf{Partie I --- Le modèle et ses moniteurs.}
\begin{enumerate}
\item Décrire le modèle de production et les trois défaillances.
\item Indiquer ce que surveille chaque couche (entrées, sorties, résultats) et à partir de quel moment chacune est disponible.
\item Définir l'\omterm{def:ml:monitoring-and-retraining:psi}{indice de stabilité de la population} (PSI) et donner les seuils calibrés.
\item Expliquer pourquoi on compte les alertes au début d'une alarme.
\end{enumerate}

\medskip\noindent\textbf{Partie II --- \omterm{def:ml:monitoring-and-retraining:performance-alarm}{Alarmes de performance}.}
\begin{enumerate}[resume]
\item Comment le CUSUM sur le coefficient d'information est-il paramétré et calibré~?
\item Combien de jours un moniteur de résultats met-il à voir chaque défaillance~?
\item Pourquoi les résultats ne peuvent-ils pas voir le changement d'unité~?
\item Quelle est la différence entre une première alarme et une alarme soutenue~?
\end{enumerate}

\medskip\noindent\textbf{Partie III --- Les délais.}
\begin{enumerate}[resume]
\item Quels moniteurs voient le changement d'unité, et à quel moment~?
\item Lesquels voient le renversement~?
\item Qu'est-ce qui voit la lente érosion, et par quel moyen~?
\item Qu'est-ce qui aurait détecté l'érosion sans son changement de fourchette~?
\end{enumerate}

\medskip\noindent\textbf{Partie IV --- Le verdict.}
\begin{enumerate}[resume]
\item Énoncer le \emph{résultat nommé}~: le délai de détection de chaque moniteur pour les trois défaillances injectées, à une fausse alarme par mois.
\item Que se passe-t-il en fantôme après le renversement, et à quelle fréquence un challenger nul est-il promu~?
\item Que permet d'économiser un canari lorsqu'un challenger défectueux est promu~?
\item Qu'aurait dû faire le desk le mardi~?
\item Quand un modèle doit-il être réentraîné, et quand doit-il être retiré~?
\item Qui doit être appelé par chaque moniteur~?
\item Que doit consigner la surveillance dans le registre~?
\item En une phrase~: de quoi un moniteur a-t-il besoin pour valoir la peine d'exister~?
\end{enumerate}
\end{problem}

\section{Questions d'entretien}

\begin{interviewq}[$\star$ \iqroles{ingénieur ML}]\label{iq:ml:monitoring-and-retraining:1}
Que surveilleriez-vous pour un modèle en production, et pourquoi à plusieurs niveaux~?
\end{interviewq}

\begin{interviewq}[$\star\star$ \iqroles{ingénieur ML, chercheur}]\label{iq:ml:monitoring-and-retraining:2}
Qu'est-ce que l'\omterm{def:ml:monitoring-and-retraining:psi}{indice de stabilité de la population}, et comment choisiriez-vous son seuil d'alarme~?
\end{interviewq}

\begin{interviewq}[$\star\star$ \iqroles{chercheur}]\label{iq:ml:monitoring-and-retraining:3}
Le CI quotidien de votre modèle est de 0.10 avec un écart type de 0.07. Au bout de combien de temps pourriez-vous affirmer qu'il est tombé à 0.08~?
\end{interviewq}

\begin{interviewq}[$\star\star$ \iqroles{ingénieur ML}]\label{iq:ml:monitoring-and-retraining:4}
Expliquez le \omterm{def:ml:monitoring-and-retraining:shadow}{déploiement fantôme}, le canari et le \omterm{def:ml:monitoring-and-retraining:shadow}{champion--challenger}. Quand chacun est-il le bon outil~?
\end{interviewq}

\begin{interviewq}[$\star\star$ \iqroles{ingénieur ML, développeur}]\label{iq:ml:monitoring-and-retraining:5}
Votre alarme de dérive des entrées se déclenche. Réentraînez-vous~?
\end{interviewq}

\begin{interviewq}[$\star\star\star$ \iqroles{ingénieur ML}]\label{iq:ml:monitoring-and-retraining:6}
Vos moniteurs appellent l'équipe vingt fois par semaine et plus personne ne regarde. Que changez-vous~?
\end{interviewq}
